% Generator template derived from books/pl/partial_differential_equations.tex.
% Keep in sync with canonical source-book preamble.
% !TeX program = pdflatex
% Wykład (rozdziały 1--32) + ramkowany kod numeryczny.
% Zadania i audyty znajdują się w oddzielnych plikach TEX/PDF.
% Kompilacja: latexmk -cd -pdf books/pl/partial_differential_equations.tex
\pdfobjcompresslevel=0
\documentclass[11pt,a4paper,oneside]{book}

% Wspólna preambuła skryptu „Równania różniczkowe cząstkowe”.
% Preambula włączona bezpośrednio do jednego pliku.
\usepackage[T1]{fontenc}
\usepackage[utf8]{inputenc}
\usepackage[polish]{babel}
\usepackage{lmodern}
\usepackage{amsmath,amssymb,amsthm,mathtools}
\usepackage{esint}
\usepackage{microtype}
\usepackage[a4paper,margin=2.7cm]{geometry}
\usepackage{enumitem}
\usepackage{xcolor}
\usepackage{graphicx}
\graphicspath{{../../assets/pl/partial_differential_equations/}}
\usepackage{tikz}
\usetikzlibrary{arrows.meta,calc,positioning}
\usepackage[unicode,colorlinks=true,linkcolor=blue!45!black,urlcolor=blue!45!black]{hyperref}

\setlength{\parindent}{0pt}
\setlength{\parskip}{4pt}
\setlength{\emergencystretch}{2em}
\setlist[enumerate]{topsep=4pt,itemsep=2pt,leftmargin=*}
\setlist[itemize]{topsep=4pt,itemsep=2pt,leftmargin=*}
\numberwithin{equation}{chapter}
\theoremstyle{plain}
\newtheorem{twierdzenie}{Twierdzenie}[chapter]
\newtheorem{stwierdzenie}[twierdzenie]{Stwierdzenie}
\newtheorem{lemat}[twierdzenie]{Lemat}
\newtheorem{wniosek}[twierdzenie]{Wniosek}
\theoremstyle{definition}
\newtheorem{definicja}[twierdzenie]{Definicja}
\newtheorem{przyklad}[twierdzenie]{Przykład}
\newtheorem{kontrprzyklad}[twierdzenie]{Kontrprzykład}
\theoremstyle{remark}
\newtheorem{uwaga}[twierdzenie]{Uwaga}

\newcommand{\R}{\mathbb{R}}
\newcommand{\C}{\mathbb{C}}
\newcommand{\N}{\mathbb{N}}
\newcommand{\T}{\mathbb{T}}
\newcommand{\eps}{\varepsilon}
\newcommand{\pd}{\partial}
\DeclareMathOperator{\supp}{supp}
\DeclareMathOperator{\tr}{tr}
\DeclareMathOperator{\diver}{div}

\title{Równania różniczkowe cząstkowe\\\large Tom I. Podstawy teorii PDE i rozwiązań słabych}
\author{}
\date{19 listopada 2026}

% BEGIN GLOBAL CODE LISTINGS (generated from latex/pl/code-listings.tex)
\usepackage{xcolor}
\usepackage{listings}
\definecolor{codebg}{RGB}{248,249,251}
\definecolor{codeborder}{RGB}{180,191,203}
\definecolor{codecomment}{RGB}{61,113,79}
\definecolor{codekeyword}{RGB}{41,78,132}
\lstdefinelanguage{Julia}{
  morekeywords={abstract,baremodule,begin,break,catch,const,continue,do,else,elseif,
    end,export,false,finally,for,function,global,if,import,in,isa,let,local,macro,
    module,mutable,primitive,quote,return,struct,true,try,type,using,where,while,
    missing,nothing,println,print},
  sensitive=true,morecomment=[l]{\#},morecomment=[n]{\#=}{=\#},
  morestring=[b]",
}
\lstdefinestyle{skrypt}{
  backgroundcolor=\color{codebg},frame=single,rulecolor=\color{codeborder},
  framesep=5pt,frameround=tttt,framerule=.5pt,
  basicstyle=\ttfamily\footnotesize,keywordstyle=\color{codekeyword}\bfseries,
  commentstyle=\color{codecomment}\itshape,stringstyle=\color{purple!65!black},
  numbers=left,numberstyle=\tiny\color{gray!75!black},numbersep=8pt,
  stepnumber=1,breaklines=true,breakatwhitespace=false,
  columns=fullflexible,keepspaces=true,showstringspaces=false,tabsize=4,
  linewidth=.90\linewidth,captionpos=b,aboveskip=11pt,belowskip=12pt,
  xleftmargin=12pt,xrightmargin=2pt,
}
\lstdefinestyle{pythonpde}{style=skrypt,language=Python}
\lstdefinestyle{python}{style=skrypt,language=Python}
\lstdefinestyle{julia}{style=skrypt,language=Julia}
\lstset{style=skrypt}
\renewcommand{\lstlistingname}{Kod}
\renewcommand{\lstlistlistingname}{Spis kodów}
% END GLOBAL CODE LISTINGS
